% Generated draft: requires mathematical and source review.
\documentclass[11pt]{article}
\usepackage[margin=25mm]{geometry}
\usepackage{fontspec}
\setmainfont{Latin Modern Roman}
\usepackage{amsmath,amssymb,mathtools,unicode-math}
\setmathfont{Latin Modern Math}
\usepackage{enumitem,needspace,xurl,hyperref}
\setlength{\parindent}{0pt}
\setlength{\parskip}{5pt}
\setlength{\emergencystretch}{3em}
\newcommand{\N}{\mathbb{N}}
\newcommand{\Z}{\mathbb{Z}}
\newcommand{\Q}{\mathbb{Q}}
\newcommand{\R}{\mathbb{R}}
\newcommand{\C}{\mathbb{C}}
\newcommand{\F}{\mathbb{F}}
\newcommand{\A}{\mathbb{A}}
\newcommand{\PP}{\mathbb{P}}
\newcommand{\E}{\mathbb{E}}
\newcommand{\eps}{\varepsilon}
\newcommand{\dd}{\,\mathrm{d}}
\begin{document}
\section*{UnsolvedMath ID: 30004692}
\textbf{Source problem code:} catalogue entry
\subsection{Definitions and mathematical statement}
A domain is a connected open set. A bounded domain \(\Omega\subset\C^n\), with \(n\geq 1\), is hyperconvex if it admits a continuous negative plurisubharmonic exhaustion \(\rho\colon\Omega\to(-\infty,0)\): for every \(c<0\), the set \(\{z\in\Omega:\rho(z)<c\}\) is relatively compact in \(\Omega\). A function is plurisubharmonic if it is upper semicontinuous and its restriction to each complex line is subharmonic wherever defined. Let \(d\lambda\) denote Lebesgue measure on \(\C^n\), and let \(dd^c\) denote the complex differential operator used to form the complex Monge--Amp\`ere operator; its normalization is not specified in the supplied record.

The supplied question asks whether the following implication holds for every bounded hyperconvex domain \(\Omega\subset\C^n\), every \(f\in C^\infty(\Omega)\) with \(f(z)>0\) for every \(z\in\Omega\), and every solution \(u\) of the indicated Dirichlet problem:
\[
\left.
\begin{aligned}
(dd^c u)^n&=f\,d\lambda &&\text{in }\Omega,\\
u&=0 &&\text{on }\partial\Omega
\end{aligned}
\right\}
\quad\Longrightarrow\quad u\in C^\infty(\Omega).
\]
Only interior smoothness is requested. The record imposes no smoothness of \(\partial\Omega\), no extension of \(f\) to the boundary, and no stated global integrability or boundedness condition on \(f\).

\emph{Unresolved formulation qualifications.} The record does not specify the admissible class of solutions, explicitly require \(u\) to be plurisubharmonic, define the weak meaning of \((dd^c u)^n\), or explain how the zero boundary condition is attained. For example, a locally bounded plurisubharmonic solution with its Monge--Amp\`ere measure interpreted in the Bedford--Taylor sense is a possible interpretation, but those hypotheses are not supplied and are not adopted here. Consequently, the displayed implication preserves the source question but requires clarification of its solution and boundary notions before it is a fully specified conjecture. The background's reference to arbitrary bounded domains is broader than the explicit hyperconvex-domain hypothesis retained above.
\subsection{Short English statement}
Does a solution of the complex Monge--Amp\`ere equation on a bounded hyperconvex domain, with zero boundary values and a strictly positive smooth interior density, necessarily become smooth throughout the interior? The intended weak solution class and meaning of the boundary values remain unspecified.
\subsection{Sources}
\begin{itemize}
\item Supplied local UnsolvedMath registry record, permanent ID \(30004692\), source problem code \texttt{catalogue entry}, titled ``Smooth Interior Regularity for Complex Monge--Amp\`ere Equations,'' in the Oberwolfach Reports Open Problems collection. No individual proposer is identified; the supplied proposed year is 2021.
\item Supplied source citation: \emph{Geometric Methods of Complex Analysis (2022)}, Oberwolfach Reports, pp.~1291--1345, DOI \texttt{10.4171/owr/2021/25}. \url{https://doi.org/10.4171/owr/2021/25}. The background describes these reports as documenting workshops and their open-problem discussions. The supplied title year is 2022, whereas the DOI contains 2021; this discrepancy has not been resolved.
\item Supplied literature source, labeled ``He, complex Monge--Amp\`ere regularity'': \url{https://arxiv.org/abs/1002.4825}. No fuller bibliographic description or precise supporting result is supplied.
\item Literature assessment checked at \texttt{2026-08-22}: the supplied curation labels the problem open and states that the interior smooth-regularity conjecture for positive smooth density on arbitrary bounded domains remains open, while existing results cover special domains and provide a priori estimates rather than the full assertion. This is supplied dated curation derived from the OWR material, not an independently verified literature assessment. Its domain description is broader than the explicit hyperconvex hypothesis of the supplied statement. The supplied curation requests independent verification of time-sensitive status and cited claims.
\end{itemize}
\end{document}
